\documentclass[a4paper]{article}
\usepackage{amsfonts}
\usepackage{amsmath}
\usepackage{amssymb}
\usepackage{graphicx}     

\begin{document}
  
\noindent \large Auteur: Michael Livchits \\
\noindent \large Studierichting: Informatica\\
\noindent \large Oefeningenreeks 6B nr. 17.7\\

\medskip

\normalsize

$\displaystyle s_n=\sum_{k=1}^{n}(2k-1)^3$\\

$= \displaystyle\sum_{k=1}^{n}(2k-1)\cdot (4k^2-4k+1)$\\

$= \displaystyle\sum_{k=1}^{n}(8k^3-12k^2+6k-1)$\\

$= 8\displaystyle\sum_{k=1}^{n}k^3-12\displaystyle\sum_{k=1}^{n}k^2+6\displaystyle\sum_{k=1}^{n}k-\displaystyle\sum_{k=1}^{n}1$\\

$= 2n^2(n+1)^2-2n(n+1)(2n+1)+3n(n+1)-n$\\

$= 2n^4 + 4n^3 + 2n^2 -4n^3 - 2n^2 -4n^2 -2n + 3n^2 + 3n -n$\\

$= 2n^4-n^2 = n^2(2n^2-1)$\\

Volgt triviaal dat het altijd $\in \mathbb{N},\forall n \geq 1$\\

\begin{thebibliography}{999}
%\bibitem{a} Referentie
\end{thebibliography}
\end{document}